\documentclass[12pt]{article}

% Modelo de relatório do Trabalho 01, baseado no template de conferências da
% Sociedade Brasileira de Computação (SBC), versão atualizada em 2017.
\usepackage{sbc-template}
\usepackage{graphicx,url}
\usepackage[utf8]{inputenc}
\usepackage[T1]{fontenc}
\usepackage{amsmath,amssymb}
\usepackage{booktabs}
\usepackage{tabularx}
\usepackage{multirow}
\usepackage{float}
\usepackage[hidelinks]{hyperref}

\sloppy
\renewcommand{\tablename}{Tabela}
\renewcommand{\figurename}{Figura}
\renewcommand{\refname}{Referências}

\title{Trabalho 01 - Busca Não Informada, Busca Informada e Busca Local}

% Substitua pelos nomes da dupla. Se o trabalho for individual, remova o segundo nome.
\author{Nome Completo do Aluno 1\inst{1}, Nome Completo do Aluno 2\inst{1}}

\address{Curso de Engenharia de Computação - CEFET-MG, campus Divinópolis\\
Disciplina de Inteligência Artificial - 2026/2
\email{email1@aluno.cefetmg.br, email2@aluno.cefetmg.br}}

\begin{document}

\maketitle

\begin{abstract}
This report presents the implementation and experimental analysis of uninformed, informed, and local search algorithms developed for Assignment 01. The experiments compare BFS, DFS, Uniform-Cost Search, Greedy Best-First Search, and A* in a weighted grid, and Hill Climbing and Random-Restart Hill Climbing in the 8-Queens problem. Replace this paragraph with a concise summary of the problem, methodology, main quantitative results, and conclusions. The abstract should not exceed approximately 10 lines.
\end{abstract}

\begin{resumo}
Este relatório apresenta a implementação e a análise experimental dos algoritmos de busca desenvolvidos no Trabalho 01. Os experimentos comparam BFS, DFS, Busca de Custo Uniforme, Busca Gulosa e A* em um ambiente ponderado, além de Hill Climbing e Random-Restart Hill Climbing no problema das 8 rainhas. Substitua este texto por um resumo objetivo contendo problema, metodologia, principais resultados quantitativos e conclusão. O resumo deve ter aproximadamente 10 linhas ou menos.
\end{resumo}

\section{Introdução}
Apresente brevemente os dois problemas tratados no trabalho e o objetivo da comparação experimental. Não transforme esta seção em uma revisão extensa dos algoritmos. O foco do relatório é a implementação realizada pelos autores e a interpretação dos resultados obtidos.

Ao final desta seção, indique de forma objetiva as perguntas que a análise pretende responder, como a relação entre número de passos e custo, o efeito das heurísticas e as limitações do Hill Climbing.

\section{Representação e implementação}
\subsection{Parte 1 - Busca não informada e informada}
Descreva a representação do estado, as estruturas utilizadas para fronteira e estados visitados, o tratamento de custos e a reconstrução do caminho. Explique também como foram implementados os critérios de desempate exigidos no enunciado.

Evite inserir grandes trechos de código. Quando necessário, use apenas pequenos fragmentos para explicar uma decisão de implementação relevante.

\subsection{Heurísticas}
Informe como foram calculadas as distâncias Manhattan e Euclidiana e explique por que elas são adequadas ao modelo de movimento adotado. A discussão sobre admissibilidade deve aparecer posteriormente, com base nos resultados e nas propriedades do problema.

\subsection{Problema das 8 rainhas}
Descreva a representação de uma configuração, a função de avaliação $h(s)$, a definição da vizinhança e os critérios de parada do Hill Climbing. Explique também como os reinícios aleatórios foram realizados e como as sementes foram controladas.

\section{Metodologia experimental}
Descreva o ambiente utilizado para os experimentos, a versão do Python e qualquer dependência relevante. Informe os mapas utilizados, as heurísticas avaliadas e a forma de obtenção das métricas.

Para as 8 rainhas, indique explicitamente que foram executadas as sementes de 0 a 99 e que o Random-Restart Hill Climbing utilizou o limite estabelecido no enunciado.

O relatório deve permitir que outra pessoa reproduza os resultados utilizando o repositório entregue. Informe neste parágrafo o endereço do repositório Git e a tag utilizada na entrega final.

\noindent\textbf{Repositório:} \url{https://...} \\
\textbf{Tag da entrega final:} \texttt{entrega-final}.

\section{Resultados}
\subsection{Parte 1 - Busca não informada e informada}
A Tabela~\ref{tab:grid} deve ser preenchida com os resultados obtidos no mapa público de teste. Se os autores realizarem experimentos adicionais, apresente-os apenas quando ajudarem a responder às questões do trabalho.

\begin{table}[H]
\centering
\caption{Resultados dos algoritmos de busca no mapa de teste.}
\label{tab:grid}
\resizebox{\textwidth}{!}{%
\begin{tabular}{llrrrrrrr}
\toprule
Algoritmo & Heurística & Passos & Custo & Gerados & Expandidos & Pico fronteira & Memória & Tempo (ms) \\
\midrule
BFS     & -          & - & - & - & - & - & - & - \\
DFS     & -          & - & - & - & - & - & - & - \\
UCS     & -          & - & - & - & - & - & - & - \\
Gulosa  & Manhattan  & - & - & - & - & - & - & - \\
Gulosa  & Euclidiana & - & - & - & - & - & - & - \\
A*      & Manhattan  & - & - & - & - & - & - & - \\
A*      & Euclidiana & - & - & - & - & - & - & - \\
\bottomrule
\end{tabular}}
\end{table}

Inclua pelo menos uma representação visual dos caminhos encontrada pelos algoritmos quando ela contribuir para a discussão. Não é necessário inserir uma figura para cada execução se isso apenas repetir informação.

\subsection{Parte 2 - Busca local e 8 rainhas}
A Tabela~\ref{tab:queens} resume os experimentos obrigatórios com as sementes de 0 a 99.

\begin{table}[H]
\centering
\caption{Resultados dos métodos de busca local no problema das 8 rainhas.}
\label{tab:queens}
\begin{tabularx}{\textwidth}{Xrr}
\toprule
Métrica & Hill Climbing & Random Restart \\
\midrule
Execuções & 100 & 100 \\
Soluções encontradas & - & - \\
Taxa de sucesso (\%) & - & - \\
Média de iterações & - & - \\
Mediana de iterações & - & - \\
Média de estados avaliados & - & - \\
Média de reinícios & N/A & - \\
Máximo de reinícios & N/A & - \\
\bottomrule
\end{tabularx}
\end{table}

Apresente ainda pelo menos um caso de falha do Hill Climbing. Informe a semente, o estado inicial, o estado final, os valores inicial e final de $h(s)$ e explique por que nenhum movimento permitido produziu melhora.

\section{Discussão}
Esta é a seção central do relatório. Não se limite a repetir os números das tabelas. Explique os resultados à luz das propriedades dos algoritmos e responda explicitamente às questões do enunciado.

A discussão deve contemplar, pelo menos:
\begin{itemize}
    \item diferença entre menor número de ações e menor custo;
    \item comportamento de BFS, DFS e UCS no mapa ponderado;
    \item desempenho e limitações da Busca Gulosa;
    \item comparação entre UCS e A* em custo e número de expansões;
    \item efeito das heurísticas Manhattan e Euclidiana;
    \item admissibilidade das heurísticas;
    \item diferença conceitual entre Busca Gulosa e Hill Climbing;
    \item motivo das falhas do Hill Climbing e efeito do Random Restart;
    \item completude e optimalidade das estratégias de busca nas condições adotadas no trabalho.
\end{itemize}

Sempre que possível, associe uma afirmação a um resultado quantitativo obtido pelos autores.

\section{Conclusão}
Sintetize os principais achados do trabalho. Indique quais estratégias foram mais adequadas em cada cenário e quais limitações foram observadas. Não introduza resultados novos nesta seção.

\section*{Contribuições dos autores}
Em trabalhos em dupla, descreva objetivamente a contribuição de cada integrante na implementação, experimentos, testes, análise e redação. Em trabalho individual, informe que todas as etapas foram realizadas pelo autor.

\section*{Declaração de uso de IA}
O detalhamento do uso de ferramentas de Inteligência Artificial deve constar no arquivo \texttt{IA\_USAGE.md} entregue com o repositório. Nesta seção, escreva apenas uma síntese, por exemplo: ``Foram utilizadas ferramentas de IA para revisão textual e esclarecimento de mensagens de erro, conforme detalhado no arquivo IA\_USAGE.md''. Caso nenhuma ferramenta tenha sido utilizada, declare isso explicitamente.

\begin{thebibliography}{99}

\bibitem[Stuart and Norvig 2021]{russell2021}
Russell, S. and Norvig, P. (2021). \textit{Artificial Intelligence: A Modern Approach}. Pearson, 4th edition.

% Adicione aqui outras referências realmente consultadas. Todo material externo usado
% na fundamentação, implementação ou discussão deve ser citado.

\end{thebibliography}

\end{document}
